%!TeX root=../emmatop.tex
\chapter[Chapter \thechapter]{}
\lettrine[lines=4,lraise=0.3]{M}{rs} Elton was first seen at church: but though devotion might be interrupted, curiosity could not be satisfied by a bride in a pew, and it must be left for the visits in form which were then to be paid, to settle whether she were very pretty indeed, or only rather pretty, or not pretty at all.


\begin{figure}[tbph]
\centering
\includegraphics[width=.9\linewidth]{32seen}
\caption{Mrs Elton was first seen at church}
\end{figure}

Emma had feelings, less of curiosity than of pride or propriety, to make her resolve on not being the last to pay her respects; and she made a point of Harriet's going with her, that the worst of the business might be gone through as soon as possible.

She could not enter the house again, could not be in the same room to which she had with such vain artifice retreated three months ago, to lace up her boot, without recollecting. A thousand vexatious thoughts would recur. Compliments, charades, and horrible blunders; and it was not to be supposed that poor Harriet should not be recollecting too; but she behaved very well, and was only rather pale and silent. The visit was of course short; and there was so much embarrassment and occupation of mind to shorten it, that Emma would not allow herself entirely to form an opinion of the lady, and on no account to give one, beyond the nothing-meaning terms of being »elegantly dressed, and very pleasing.«

She did not really like her. She would not be in a hurry to find fault, but she suspected that there was no elegance;—ease, but not elegance.— She was almost sure that for a young woman, a stranger, a bride, there was too much ease. Her person was rather good; her face not unpretty; but neither feature, nor air, nor voice, nor manner, were elegant. Emma thought at least it would turn out so.

As for Mr Elton, his manners did not appear—but no, she would not permit a hasty or a witty word from herself about his manners. It was an awkward ceremony at any time to be receiving wedding visits, and a man had need be all grace to acquit himself well through it. The woman was better off; she might have the assistance of fine clothes, and the privilege of bashfulness, but the man had only his own good sense to depend on; and when she considered how peculiarly unlucky poor Mr Elton was in being in the same room at once with the woman he had just married, the woman he had wanted to marry, and the woman whom he had been expected to marry, she must allow him to have the right to look as little wise, and to be as much affectedly, and as little really easy as could be.

»Well, Miss Woodhouse,« said Harriet, when they had quitted the house, and after waiting in vain for her friend to begin; »Well, Miss Woodhouse, (with a gentle sigh,) what do you think of her?—Is not she very charming?«

There was a little hesitation in Emma's answer.

»Oh! yes—very—a very pleasing young woman.«

»I think her beautiful, quite beautiful.«

»Very nicely dressed, indeed; a remarkably elegant gown.«

»I am not at all surprized that he should have fallen in love.«

»Oh! no—there is nothing to surprize one at all.—A pretty fortune; and she came in his way.«

»I dare say,« returned Harriet, sighing again, »I dare say she was very much attached to him.«

»Perhaps she might; but it is not every man's fate to marry the woman who loves him best. Miss Hawkins perhaps wanted a home, and thought this the best offer she was likely to have.«

»Yes,« said Harriet earnestly, »and well she might, nobody could ever have a better. Well, I wish them happy with all my heart. And now, Miss Woodhouse, I do not think I shall mind seeing them again. He is just as superior as ever;—but being married, you know, it is quite a different thing. No, indeed, Miss Woodhouse, you need not be afraid; I can sit and admire him now without any great misery. To know that he has not thrown himself away, is such a comfort!—She does seem a charming young woman, just what he deserves. Happy creature! He called her »Augusta.« How delightful!«

When the visit was returned, Emma made up her mind. She could then see more and judge better. From Harriet's happening not to be at Hartfield, and her father's being present to engage Mr Elton, she had a quarter of an hour of the lady's conversation to herself, and could composedly attend to her; and the quarter of an hour quite convinced her that Mrs Elton was a vain woman, extremely well satisfied with herself, and thinking much of her own importance; that she meant to shine and be very superior, but with manners which had been formed in a bad school, pert and familiar; that all her notions were drawn from one set of people, and one style of living; that if not foolish she was ignorant, and that her society would certainly do Mr Elton no good.

Harriet would have been a better match. If not wise or refined herself, she would have connected him with those who were; but Miss Hawkins, it might be fairly supposed from her easy conceit, had been the best of her own set. The rich brother-in-law near Bristol was the pride of the alliance, and his place and his carriages were the pride of him.

The very first subject after being seated was Maple Grove, »My brother Mr Suckling's seat;«—a comparison of Hartfield to Maple Grove. The grounds of Hartfield were small, but neat and pretty; and the house was modern and well-built. Mrs Elton seemed most favourably impressed by the size of the room, the entrance, and all that she could see or imagine. »Very like Maple Grove indeed!—She was quite struck by the likeness!—That room was the very shape and size of the morning-room at Maple Grove; her sister's favourite room.«—Mr Elton was appealed to.—»Was not it astonishingly like?—She could really almost fancy herself at Maple Grove.«

»And the staircase—You know, as I came in, I observed how very like the staircase was; placed exactly in the same part of the house. I really could not help exclaiming! I assure you, Miss Woodhouse, it is very delightful to me, to be reminded of a place I am so extremely partial to as Maple Grove. I have spent so many happy months there! (with a little sigh of sentiment). A charming place, undoubtedly. Every body who sees it is struck by its beauty; but to me, it has been quite a home. Whenever you are transplanted, like me, Miss Woodhouse, you will understand how very delightful it is to meet with any thing at all like what one has left behind. I always say this is quite one of the evils of matrimony.«

Emma made as slight a reply as she could; but it was fully sufficient for Mrs Elton, who only wanted to be talking herself.

»So extremely like Maple Grove! And it is not merely the house—the grounds, I assure you, as far as I could observe, are strikingly like. The laurels at Maple Grove are in the same profusion as here, and stand very much in the same way—just across the lawn; and I had a glimpse of a fine large tree, with a bench round it, which put me so exactly in mind! My brother and sister will be enchanted with this place. People who have extensive grounds themselves are always pleased with any thing in the same style.«

Emma doubted the truth of this sentiment. She had a great idea that people who had extensive grounds themselves cared very little for the extensive grounds of any body else; but it was not worth while to attack an error so double-dyed, and therefore only said in reply,

»When you have seen more of this country, I am afraid you will think you have overrated Hartfield. Surry is full of beauties.«

»Oh! yes, I am quite aware of that. It is the garden of England, you know. Surry is the garden of England.«

»Yes; but we must not rest our claims on that distinction. Many counties, I believe, are called the garden of England, as well as Surry.«

»No, I fancy not,« replied Mrs Elton, with a most satisfied smile. »I never heard any county but Surry called so.«

Emma was silenced.

»My brother and sister have promised us a visit in the spring, or summer at farthest,« continued Mrs Elton; »and that will be our time for exploring. While they are with us, we shall explore a great deal, I dare say. They will have their barouche-landau, of course, which holds four perfectly; and therefore, without saying any thing of our carriage, we should be able to explore the different beauties extremely well. They would hardly come in their chaise, I think, at that season of the year. Indeed, when the time draws on, I shall decidedly recommend their bringing the barouche-landau; it will be so very much preferable. When people come into a beautiful country of this sort, you know, Miss Woodhouse, one naturally wishes them to see as much as possible; and Mr Suckling is extremely fond of exploring. We explored to King's-Weston twice last summer, in that way, most delightfully, just after their first having the barouche-landau. You have many parties of that kind here, I suppose, Miss Woodhouse, every summer?«

»No; not immediately here. We are rather out of distance of the very striking beauties which attract the sort of parties you speak of; and we are a very quiet set of people, I believe; more disposed to stay at home than engage in schemes of pleasure.«

»Ah! there is nothing like staying at home for real comfort. Nobody can be more devoted to home than I am. I was quite a proverb for it at Maple Grove. Many a time has Selina said, when she has been going to Bristol, »I really cannot get this girl to move from the house. I absolutely must go in by myself, though I hate being stuck up in the barouche-landau without a companion; but Augusta, I believe, with her own good-will, would never stir beyond the park paling.« Many a time has she said so; and yet I am no advocate for entire seclusion. I think, on the contrary, when people shut themselves up entirely from society, it is a very bad thing; and that it is much more advisable to mix in the world in a proper degree, without living in it either too much or too little. I perfectly understand your situation, however, Miss Woodhouse—(looking towards Mr Woodhouse), Your father's state of health must be a great drawback. Why does not he try Bath?—Indeed he should. Let me recommend Bath to you. I assure you I have no doubt of its doing Mr Woodhouse good.«

»My father tried it more than once, formerly; but without receiving any benefit; and Mr Perry, whose name, I dare say, is not unknown to you, does not conceive it would be at all more likely to be useful now.«

»Ah! that's a great pity; for I assure you, Miss Woodhouse, where the waters do agree, it is quite wonderful the relief they give. In my Bath life, I have seen such instances of it! And it is so cheerful a place, that it could not fail of being of use to Mr Woodhouse's spirits, which, I understand, are sometimes much depressed. And as to its recommendations to you, I fancy I need not take much pains to dwell on them. The advantages of Bath to the young are pretty generally understood. It would be a charming introduction for you, who have lived so secluded a life; and I could immediately secure you some of the best society in the place. A line from me would bring you a little host of acquaintance; and my particular friend, Mrs Partridge, the lady I have always resided with when in Bath, would be most happy to shew you any attentions, and would be the very person for you to go into public with.«

\begin{figure}[tbph]
\centering
\includegraphics[width=.9\linewidth]{32widow}
\caption{Some vulgar, dashing widow}
\end{figure}

It was as much as Emma could bear, without being impolite. The idea of her being indebted to Mrs Elton for what was called an introduction—of her going into public under the auspices of a friend of Mrs Elton's—probably some vulgar, dashing widow, who, with the help of a boarder, just made a shift to live!—The dignity of Miss Woodhouse, of Hartfield, was sunk indeed!

She restrained herself, however, from any of the reproofs she could have given, and only thanked Mrs Elton coolly; »but their going to Bath was quite out of the question; and she was not perfectly convinced that the place might suit her better than her father.« And then, to prevent farther outrage and indignation, changed the subject directly.

»I do not ask whether you are musical, Mrs Elton. Upon these occasions, a lady's character generally precedes her; and Highbury has long known that you are a superior performer.«

»Oh! no, indeed; I must protest against any such idea. A superior performer!—very far from it, I assure you. Consider from how partial a quarter your information came. I am doatingly fond of music—passionately fond;—and my friends say I am not entirely devoid of taste; but as to any thing else, upon my honour my performance is mediocre to the last degree. You, Miss Woodhouse, I well know, play delightfully. I assure you it has been the greatest satisfaction, comfort, and delight to me, to hear what a musical society I am got into. I absolutely cannot do without music. It is a necessary of life to me; and having always been used to a very musical society, both at Maple Grove and in Bath, it would have been a most serious sacrifice. I honestly said as much to Mr E. when he was speaking of my future home, and expressing his fears lest the retirement of it should be disagreeable; and the inferiority of the house too—knowing what I had been accustomed to—of course he was not wholly without apprehension. When he was speaking of it in that way, I honestly said that the world I could give up—parties, balls, plays—for I had no fear of retirement. Blessed with so many resources within myself, the world was not necessary to me. I could do very well without it. To those who had no resources it was a different thing; but my resources made me quite independent. And as to smaller-sized rooms than I had been used to, I really could not give it a thought. I hoped I was perfectly equal to any sacrifice of that description. Certainly I had been accustomed to every luxury at Maple Grove; but I did assure him that two carriages were not necessary to my happiness, nor were spacious apartments. »But,« said I, »to be quite honest, I do not think I can live without something of a musical society. I condition for nothing else; but without music, life would be a blank to me.««

»We cannot suppose,« said Emma, smiling, »that Mr Elton would hesitate to assure you of there being a very musical society in Highbury; and I hope you will not find he has outstepped the truth more than may be pardoned, in consideration of the motive.«

»No, indeed, I have no doubts at all on that head. I am delighted to find myself in such a circle. I hope we shall have many sweet little concerts together. I think, Miss Woodhouse, you and I must establish a musical club, and have regular weekly meetings at your house, or ours. Will not it be a good plan? If we exert ourselves, I think we shall not be long in want of allies. Something of that nature would be particularly desirable for me, as an inducement to keep me in practice; for married women, you know—there is a sad story against them, in general. They are but too apt to give up music.«

»But you, who are so extremely fond of it—there can be no danger, surely?«

»I should hope not; but really when I look around among my acquaintance, I tremble. Selina has entirely given up music—never touches the instrument—though she played sweetly. And the same may be said of Mrs Jeffereys—Clara Partridge, that was—and of the two Milmans, now Mrs Bird and Mrs James Cooper; and of more than I can enumerate. Upon my word it is enough to put one in a fright. I used to be quite angry with Selina; but really I begin now to comprehend that a married woman has many things to call her attention. I believe I was half an hour this morning shut up with my housekeeper.«

»But every thing of that kind,« said Emma, »will soon be in so regular a train\longdash«

»Well,« said Mrs Elton, laughing, »we shall see.«

Emma, finding her so determined upon neglecting her music, had nothing more to say; and, after a moment's pause, Mrs Elton chose another subject.

»We have been calling at Randalls,« said she, »and found them both at home; and very pleasant people they seem to be. I like them extremely. Mr Weston seems an excellent creature—quite a first-rate favourite with me already, I assure you. And she appears so truly good—there is something so motherly and kind-hearted about her, that it wins upon one directly. She was your governess, I think?«

Emma was almost too much astonished to answer; but Mrs Elton hardly waited for the affirmative before she went on.

»Having understood as much, I was rather astonished to find her so very lady-like! But she is really quite the gentlewoman.«

»Mrs Weston's manners,« said Emma, »were always particularly good. Their propriety, simplicity, and elegance, would make them the safest model for any young woman.«

»And who do you think came in while we were there?«

Emma was quite at a loss. The tone implied some old acquaintance—and how could she possibly guess?

»Knightley!« continued Mrs Elton; »Knightley himself!—Was not it lucky?—for, not being within when he called the other day, I had never seen him before; and of course, as so particular a friend of Mr E.'s, I had a great curiosity. »My friend Knightley« had been so often mentioned, that I was really impatient to see him; and I must do my caro sposo the justice to say that he need not be ashamed of his friend. Knightley is quite the gentleman. I like him very much. Decidedly, I think, a very gentleman-like man.«

Happily, it was now time to be gone. They were off; and Emma could breathe.

»Insufferable woman!« was her immediate exclamation. »Worse than I had supposed. Absolutely insufferable! Knightley!—I could not have believed it. Knightley!—never seen him in her life before, and call him Knightley!—and discover that he is a gentleman! A little upstart, vulgar being, with her Mr E., and her caro sposo, and her resources, and all her airs of pert pretension and underbred finery. Actually to discover that Mr Knightley is a gentleman! I doubt whether he will return the compliment, and discover her to be a lady. I could not have believed it! And to propose that she and I should unite to form a musical club! One would fancy we were bosom friends! And Mrs Weston!—Astonished that the person who had brought me up should be a gentlewoman! Worse and worse. I never met with her equal. Much beyond my hopes. Harriet is disgraced by any comparison. Oh! what would Frank Churchill say to her, if he were here? How angry and how diverted he would be! Ah! there I am—thinking of him directly. Always the first person to be thought of! How I catch myself out! Frank Churchill comes as regularly into my mind!«—

All this ran so glibly through her thoughts, that by the time her father had arranged himself, after the bustle of the Eltons' departure, and was ready to speak, she was very tolerably capable of attending.

»Well, my dear,« he deliberately began, »considering we never saw her before, she seems a very pretty sort of young lady; and I dare say she was very much pleased with you. She speaks a little too quick. A little quickness of voice there is which rather hurts the ear. But I believe I am nice; I do not like strange voices; and nobody speaks like you and poor Miss Taylor. However, she seems a very obliging, pretty-behaved young lady, and no doubt will make him a very good wife. Though I think he had better not have married. I made the best excuses I could for not having been able to wait on him and Mrs Elton on this happy occasion; I said that I hoped I should in the course of the summer. But I ought to have gone before. Not to wait upon a bride is very remiss. Ah! it shews what a sad invalid I am! But I do not like the corner into Vicarage Lane.«

»I dare say your apologies were accepted, sir. Mr Elton knows you.«

»Yes: but a young lady—a bride—I ought to have paid my respects to her if possible. It was being very deficient.«

»But, my dear papa, you are no friend to matrimony; and therefore why should you be so anxious to pay your respects to a bride? It ought to be no recommendation to you. It is encouraging people to marry if you make so much of them.«

»No, my dear, I never encouraged any body to marry, but I would always wish to pay every proper attention to a lady—and a bride, especially, is never to be neglected. More is avowedly due to her. A bride, you know, my dear, is always the first in company, let the others be who they may.«

»Well, papa, if this is not encouragement to marry, I do not know what is. And I should never have expected you to be lending your sanction to such vanity-baits for poor young ladies.«

»My dear, you do not understand me. This is a matter of mere common politeness and good-breeding, and has nothing to do with any encouragement to people to marry.«

Emma had done. Her father was growing nervous, and could not understand her. Her mind returned to Mrs Elton's offences, and long, very long, did they occupy her.